\chapter{Results}

\section{Simulation Results}
\subsection{Simulation Setup}

\subsection{Cloud-based Planning Pipeline}

For the planner pipeline, both cloud-based and locally deployable models were evaluated. 
Due to hardware constraints and real-time execution requirements, 
small quantized LLMs with 1--3 billion parameters that could run locally on a PC proved too slow and unreliable, 
even for tasks consisting of only one or two subtasks. 
Therefore, the planner was evaluated using GPT-5, GPT-5-mini, GPT-4o, and GPT-4.1 through the OpenAI API \cite{openai_models}. 
The models were tested on tasks with the number of subtasks ranging from 1 to 10. 
For each number of subtasks, three tasks were tested.

Task descriptions were phrased in unstructured natural language with varied wording and sentence structures to leverage
the reasoning capabilities of the LLMs. 
Examples ranged from simple instructions, e.g., ”Acquire a thermal image of House3.”, to more complex tasks with several subtasks, 
e.g., ”Record videos for all solar panels and rooftops, take thermal imagery of each, and inspect the Tower and Base.”.
For comparison, an optimal schedule for each task was computed using a VRP solver, which computed optimal solutions.
This baseline was used to guide model evaluation and selection.

\Cref{tab:avg_vrp_gap_by_subtasks} shows the schedule quality of each model by comparing the generated mission
makespan to a rule-based VRP planner \cite{google_vrp}.
The table also reports the number of successfully generated schedules out of three for each subtask count.
\Cref{tab:avg_inference_time_by_subtasks} shows the average inference time computed over the successfully generated schedules for each model.

\begin{table}[htbp]
\centering
\caption[Average makespan difference between the planning pipeline and the VRP baseline by number of subtasks.]
{Average makespan difference between the planning pipeline and the VRP baseline by number of subtasks. 
Values in parentheses indicate the number of solved cases out of three tasks.
Averages are computed over valid schedules only. Cases with no valid solutions are indicated by --.}
\label{tab:avg_vrp_gap_by_subtasks}
\centering
\begin{tabular}{lllll}
\toprule
model & GPT-5 & GPT-5-mini & GPT-4o & GPT-4.1 \\
Subtasks &  &  &  &  \\
\midrule
1 & 0.00 (3/3) & 0.00 (3/3) & 0.00 (3/3) & 0.00 (3/3) \\
2 & 0.00 (3/3) & 0.00 (3/3) & 14.07 (3/3) & 0.00 (3/3) \\
3 & 0.00 (3/3) & 0.00 (3/3) & 36.54 (3/3) & 62.60 (3/3) \\
4 & 0.00 (3/3) & 0.00 (3/3) & 121.29 (3/3) & 121.29 (3/3) \\
5 & 0.00 (3/3) & 0.00 (3/3) & 132.79 (3/3) & 112.17 (3/3) \\
6 & 0.00 (3/3) & 0.86 (3/3) & 78.51 (3/3) & --  \\
7 & 0.00 (3/3) & 5.77 (2/3) & 62.95 (3/3) & 52.03 (2/3) \\
8 & 0.00 (3/3) & 14.90 (3/3) & --  & --  \\
9 & 0.56 (3/3) & 15.88 (2/3) & 22.16 (2/3) & 79.66 (1/3) \\
10 & 4.35 (3/3) & 25.19 (3/3) & --  & 161.43 (1/3) \\
\bottomrule
\end{tabular}

\end{table}

\begin{table}[htbp]
\centering
\caption[Average inference time of the planning pipeline by number of subtasks.]
{Average inference time of the planning pipeline by number of subtasks. Values in parentheses indicate solved cases over three tasks.
Averages are computed over valid schedules only. Cases with no valid solutions are indicated by --.}
\label{tab:avg_inference_time_by_subtasks}
\centering
\begin{tabular}{lllll}
\toprule
Model & GPT-5 & GPT-5-mini & GPT-4o & GPT-4.1 \\
Subtasks &  &  &  &  \\
\midrule
1 & 32.97 (3/3) & 20.27 (3/3) & 6.17 (3/3) & 4.07 (3/3) \\
2 & 49.90 (3/3) & 21.87 (3/3) & 5.00 (3/3) & 4.03 (3/3) \\
3 & 94.50 (3/3) & 33.50 (3/3) & 8.93 (3/3) & 7.13 (3/3) \\
4 & 111.90 (3/3) & 51.63 (3/3) & 12.90 (3/3) & 7.13 (3/3) \\
5 & 112.33 (3/3) & 45.83 (3/3) & 15.63 (3/3) & 11.30 (3/3) \\
6 & 205.07 (3/3) & 98.10 (3/3) & 14.47 (3/3) & -- \\
7 & 284.37 (3/3) & 119.30 (2/3) & 17.90 (3/3) & 17.00 (2/3) \\
8 & 284.77 (3/3) & 133.67 (3/3) & --  & -- \\
9 & 243.07 (3/3) & 112.05 (2/3) & 19.70 (2/3) & 16.50 (1/3) \\
10 & 269.00 (3/3) & 97.73 (3/3) & -- & 16.80 (1/3) \\
\bottomrule
\end{tabular}

\end{table}


The results show that GPT-5 consistently generated schedules close to the VRP baseline, 
but its inference time was significantly higher than that of GPT-5-mini. 
This makes GPT-5 less suitable for interactive planning and replanning. 
GPT-5-mini produced valid schedules for most tasks and achieved near-optimal schedules for tasks with up to five subtasks, 
while requiring considerably less inference time. 
The smaller and older models, GPT-4o and GPT-4.1, 
performed well on simpler tasks but increasingly produced suboptimal or invalid schedules as the number of subtasks increased. 
Based on this trade-off between schedule quality, inference time, and API cost, GPT-5-mini was selected for the cloud-based planning pipeline.

\subsection{Onboard Task Admission}
The task admission module has different requirements from the cloud-based planner. 
It must evaluate drone telemetry quickly and reliably during execution, and therefore a smaller locally deployable model is preferred. 
The evaluated models were tested on four categories: acceptable tasks, drone-state errors, link-quality errors, and flight-time errors. 
Each category contained 20 test cases.
Seven 4-bit quantized models, ranging from 0.6 to 4 billion parameters, were evaluated from the Qwen \cite{ollama_qwen25,ollama_qwen3}, 
Gemma \cite{ollama_gemma3}, Llama \cite{ollama_llama32}, and Phi \cite{ollama_phi4} model families.

\Cref{tab:task_admission_accuracy} reports the number of correctly classified cases out of 20 for each category. 
\Cref{tab:task_admission_inference_time} reports the average and maximum inference time for each model.

\begin{table}[t]
\centering
\caption[Task admission accuracy across failure categories.]
{Task admission accuracy across failure categories. Each value indicates the number of correctly solved cases out of 20.}
\label{tab:task_admission_accuracy}
\small
\begin{tabular}{lrrrr}
\toprule
Model 
& \makecell[r]{Acceptable\\task} 
& \makecell[r]{Drone state\\error} 
& \makecell[r]{Link quality\\error} 
& \makecell[r]{Flight time\\error} \\
\midrule
qwen3:0.6b       & 18 & 8  & 6  & 16 \\
gemma3:1b        & 20 & 15 & 0  & 0  \\
qwen3:1.7b       & 20 & 20 & 20 & 20 \\
qwen2.5:3b       & 7  & 20 & 20 & 9  \\
llama3.2:3b      & 13 & 20 & 18 & 0  \\
phi4-mini:3.8b   & 20 & 20 & 20 & 1  \\
gemma3:4b        & 14 & 20 & 9  & 0  \\
\bottomrule
\end{tabular}
\end{table}
\begin{table}[t]
\centering
\caption{Average and maximum inference time for the task admission module.}
\label{tab:task_admission_inference_time}
\begin{tabular}{lrr}
\toprule
Model & Average inference time (s) & Maximum inference time (s) \\
\midrule
qwen3:0.6b & 7.99 & 24.75 \\
gemma3:1b & 0.76 & 2.83 \\
qwen3:1.7b & 3.60 & 7.45 \\
qwen2.5:3b & 1.02 & 6.44 \\
llama3.2:3b & 1.25 & 14.24 \\
phi4-mini:3.8b & 1.52 & 7.72 \\
gemma3:4b & 1.20 & 9.12 \\
\bottomrule
\end{tabular}

\end{table}


The results show that qwen3:1.7b achieved perfect classification accuracy across all four categories, correctly solving all 80 test cases. 
Although some models produced lower average inference times, they failed on important failure categories such as link-quality or flight-time errors. 
Since the task admission module is safety-relevant and directly affects replanning decisions, 
classification reliability was prioritized over minimal inference time. 
Therefore, qwen3:1.7b was selected for onboard task admission.


\section{Experimental Results}
\subsection{Drones}
\subsection{Experimental setup}